% Tableau des labs du cours NRAC dans NDAP SUKULU (propre au cours).
Le cours compte dix laboratoires, disponibles sous le code \codeUE. Le dernier
rassemble le dispositif de GRC de la BANQUE SANAGA et sert de synthèse.

\begin{tableaucours}{lL{6.2cm}L{1.6cm}L{2.0cm}}{%
    \tableauHeader{Lab} &
    \tableauHeader{Sujet} &
    \tableauHeader{Chapitre} &
    \tableauHeader{Opération}
}
NRAC-LAB-01 & Coter le registre des risques et sa base de preuve & 2 & SANAGA \\
NRAC-LAB-02 & EBIOS Risk Manager~: cadrage et sources de risque & 4 & LIZBIZ \\
NRAC-LAB-03 & EBIOS Risk Manager~: écosystème et scénario opérationnel & 4 & LIZBIZ \\
NRAC-LAB-04 & Quantifier un risque cyber avec FAIR & 5 & SANAGA \\
NRAC-LAB-05 & Décomposer, coter, conclure & 7 & SANAGA \\
NRAC-LAB-06 & Confronter une déclaration d'applicabilité aux preuves & 8 & SANAGA \\
NRAC-LAB-07 & Une base commune, plusieurs référentiels & 9 & ANPP \\
NRAC-LAB-08 & Auditer un SMSI & 11 & SANAGA \\
NRAC-LAB-09 & Audit technique sur exports figés & 12 & LIZBIZ \\
NRAC-LAB-10 & Synthèse GRC~: obligations, matrice, indicateurs & 13 & SANAGA \\
\end{tableaucours}
\vspace{-0.8em}
\begin{center}{\small\itshape Les laboratoires du cours dans NDAP SUKULU}\end{center}
